% META: {"id": "TEXP-ARITHMETIQUE-RE-C3", "chapitre": "TEXP-ARITHMETIQUE", "type_objet": "remediation", "capacites_codes": ["C3"], "status": "approved"}

%---------------------------------------------------------------
% FICHE DE REMEDIATION -- C3 : Erreurs sur les inegalites par convexite
%---------------------------------------------------------------

\subsection*{Rappel -- Erreurs sur les inegalites par convexite}

\textbf{Erreur 1 :} Confondre convexite et concavite. Si $f'' > 0$, $f$ est convexe (en forme de $\cup$).

\textbf{Erreur 2 :} Appliquer Jensen dans le mauvais sens pour une fonction concave.

\textbf{Erreur 3 :} Oublier de verifier la convexite avant d'utiliser l'inegalite des tangentes.

\begin{exercice}{TEXP-ARI-RE-C3-EX1}{1}{12}
\begin{enumerate}
  \item Montrer que $f(x) = x^2 + \mathrm{e}^x$ est convexe sur $\mathbb{R}$.
  \item En deduire que $f\!\left(\dfrac{a+b}{2}\right) \leq \dfrac{f(a) + f(b)}{2}$.
  \item Verifier pour $a = 0$ et $b = 1$.
\end{enumerate}
\end{exercice}

% BEGIN-VERIFY
% from sympy import *
% x = symbols('x')
% f = x**2 + exp(x)
% fpp = diff(f, x, 2)
% assert fpp == 2 + exp(x)
% assert fpp.subs(x, 0) == 3  # > 0, convexe
% # Jensen: f(1/2) <= (f(0) + f(1))/2
% lhs = f.subs(x, Rational(1,2))
% rhs = (f.subs(x, 0) + f.subs(x, 1))/2
% assert simplify(rhs - lhs) > 0
% END-VERIFY
